% !TeX root = ../regression_logistique.tex
\chapter{Interprétation du modèle ajusté}
\label{ch:modele}

\orient{%
lire le signe et la valeur d'un coefficient standardisé sans surinterpréter ;
énoncer le contenu complet d'un modèle enregistré ; distinguer paramètres ajustés
et décisions préalables}{%
cote et logit (\annexeref{ann:prerequis}) ; chapitre~\ref{ch:descente}}{%
essentiel}

\section{Interprétation des coefficients}

La descente s'arrête sur trois nombres, qui portent tout ce que l'ajustement a
extrait des $1\,600$ observations :
\[
  w_0=-4{,}745,\qquad
  w_1=-3{,}454\ \text{(\Herb)},\qquad
  w_2=+3{,}469\ \text{(\Runes)} .
\]

Le \emph{signe} donne le sens de l'effet. Le coefficient d'\Herb est négatif :
à note d'\Runes fixée, une note d'\Herb plus haute abaisse le score, donc la
probabilité annoncée. Celui d'\Runes est positif : l'effet est inverse. Dans le
nuage, les \Gryf occupent ainsi la région en haut à gauche.

La \emph{valeur absolue} mesure l'effet d'un écart-type de cette variable sur la
log-cote, \emph{les autres variables restant fixées}. Les deux valeurs $3{,}45$
et $3{,}47$ sont presque égales : un écart-type de l'une ou de l'autre déplace le
score de la même quantité. Cette comparaison n'a de sens que parce que les notes
ont été standardisées : sur les notes brutes, le coefficient d'\Runes serait
vingt fois plus petit que celui d'\Herb pour la seule raison que ses valeurs sont
vingt fois plus dispersées.

La portée de cette lecture est celle d'un effet conditionnel. Avec des variables
corrélées, retirer une colonne modifie les coefficients de celles qui restent,
parfois jusqu'au changement de signe ; un coefficient nul signifie « aucun apport
une fois les autres variables connues ». L'ajustement porte par ailleurs sur des
associations observées, et laisse ouverte la question de savoir si agir sur une
note modifierait la classe. Classer des variables par $|w_j|$ suppose donc des
colonnes standardisées et peu corrélées, et reste une lecture descriptive du
modèle.

Le \emph{biais} $w_0$ est le seul des trois à s'appliquer identiquement à toutes
les observations. Il vaut la log-cote prédite pour une observation dont toutes
les notes standardisées sont nulles, c'est-à-dire moyennes dans chaque matière.

Le biais égale le logit de la proportion de \Gryf dans le seul cas d'un modèle
réduit au biais, sans variable --- le cas du premier pas de la descente. Dès que
des variables corrélées à la classe entrent dans
le score, il s'en écarte : ici $\sigma(-4{,}745)=0{,}0087$, contre une proportion
observée de $0{,}2044$.

Ces trois nombres se lisent aussi de façon quantitative, en revenant à la cote.
Le score étant le logarithme de la cote, ajouter une unité à une note
standardisée augmente le score de $w_j$, donc multiplie la cote par $e^{w_j}$ :
\[
  \frac{p'/(1-p')}{p/(1-p)}=e^{\,w_j}
  \qquad\text{lorsque } x_j \text{ augmente de } 1 .
\]
Un écart-type de plus en \Runes multiplie ainsi par $e^{3{,}469}=32$ la cote
d'appartenir à \Gryf, la note d'\Herb restant fixée ; un écart-type de plus en
\Herb la divise par $e^{3{,}454}=32$ également, les deux coefficients étant
presque opposés. Le biais donne la cote prédite pour une observation moyenne dans
les deux matières : $e^{-4{,}745}=0{,}0087$, soit une chance sur cent quinze.

Les deux panneaux ci-dessous font varier $w_0$ puis $w_1$ autour du modèle
obtenu, les autres coefficients restant à leur valeur. Un troisième panneau pour
$w_2$ montrerait un pivotement identique à celui de $w_1$ : c'est le rapport
$w_1/w_2$ qui fixe la direction de la frontière, et le faire varier par l'un ou
par l'autre revient au même.

\begin{figure}[htbp]
\centering
\begin{graphique}{Deux nuages de points identiques, chacun traversé par trois droites. À gauche, seul le biais varie : les trois droites sont parallèles. À droite, seul le coefficient d'Herbology varie : les trois droites pivotent autour d'un même point. Le trait épais est le modèle obtenu.}
\begin{tikzpicture}[baseline=(current bounding box.north),
  alt={Nuage de points traversé par trois droites parallèles. Seul le biais varie ; le trait épais représente le modèle ajusté.}]
\begin{axis}[panel,width=.46\textwidth,height=6cm,
  title={variation de $w_0$\\[-1pt] \scriptsize
    $w_1=-3{,}454$,\ $w_2=+3{,}469$},
  title style={align=center},
  xlabel={\Herb},ylabel={\Runes}]
  \addplot[only marks,mark=x,mark size=1.1pt,line width=.3pt,warning!80,opacity=.45]
    table[x=x,y=y] {data/autres.dat};
  \addplot[only marks,mark=*,mark size=.8pt,accent,opacity=.6]
    table[x=x,y=y] {data/gryffindor.dat};
  \addplot[ink!35,thick,domain=-2.8:2.8]{(2.0+3.4535*x)/3.4688};
  \addplot[ink,very thick,domain=-2.8:2.8]{(4.7454+3.4535*x)/3.4688};
  \addplot[ink!35,thick,domain=-2.8:2.8]{(7.5+3.4535*x)/3.4688};
\end{axis}
\end{tikzpicture}\hfill
\begin{tikzpicture}[baseline=(current bounding box.north),
  alt={Nuage de points traversé par trois droites qui pivotent autour d'un même point. Seul le coefficient d'Herbology varie ; le trait épais représente le modèle ajusté.}]
\begin{axis}[panel,width=.46\textwidth,height=6cm,
  title={variation de $w_1$\\[-1pt] \scriptsize
    $w_0=-4{,}745$,\ $w_2=+3{,}469$},
  title style={align=center},
  xlabel={\Herb},yticklabels={}]
  \addplot[only marks,mark=x,mark size=1.1pt,line width=.3pt,warning!80,opacity=.45]
    table[x=x,y=y] {data/autres.dat};
  \addplot[only marks,mark=*,mark size=.8pt,accent,opacity=.6]
    table[x=x,y=y] {data/gryffindor.dat};
  \addplot[ink!35,thick,domain=-2.8:2.8]{(4.7454+1.2*x)/3.4688};
  \addplot[ink,very thick,domain=-2.8:2.8]{(4.7454+3.4535*x)/3.4688};
  \addplot[ink!35,thick,domain=-2.8:2.8]{(4.7454+7.0*x)/3.4688};
\end{axis}
\end{tikzpicture}
\end{graphique}
\caption{Variation isolée du biais, à gauche : translation. Variation isolée de
$w_1$, à droite : rotation et modification de la pente de la sigmoïde.}
\label{fig:roles}
\end{figure}
\FloatBarrier

\begin{figure}[htbp]
\centering
\begin{graphique}{Trois blocs. En haut à gauche, les trois coefficients ajustés. En haut à droite, les conventions : ordre des matières et réponse associée. En bas, un tableau des statistiques d'entraînement donnant médiane, moyenne et écart-type pour chacune des deux matières.}
\begin{tikzpicture}[
  bx/.style={draw=ink!30,rounded corners=2.5pt,inner sep=6pt,align=center},
  ttl/.style={font=\footnotesize\bfseries,text=ink,anchor=base}]

\node[bx,fill=accentlight,minimum height=15mm] (w) at (0,0)
  {\footnotesize$\begin{array}{r@{\;=\;}r}
     w_0 & -4{,}745\\ w_1 & -3{,}454\\ w_2 & +3{,}469
   \end{array}$};
\node[bx,fill=warninglight,minimum height=15mm,right=8mm of w] (cv)
  {\footnotesize\begin{tabular}{@{}ll@{}}
     ordre des matières : & \Herb, \Runes\\[2pt]
     réponse associée : & \Gryf ou non
   \end{tabular}};
\node[ttl,above=2.5mm of w] {coefficients ajustés};
\node[ttl,above=2.5mm of cv] {conventions};

\coordinate (mid) at ($(w.south)!0.5!(cv.south)$);
\node[bx,fill=rulelight,below=10mm of mid,anchor=north] (st)
  {\footnotesize\begin{tabular}{@{}lrrr@{}}
     \toprule
     matière & médiane & moyenne & écart-type\\
     \midrule
     \Herb  & $3{,}469$   & $1{,}189$   & $5{,}175$\\
     \Runes & $463{,}918$ & $495{,}052$ & $105{,}186$\\
     \bottomrule
   \end{tabular}};
\node[ttl,above=2mm of st] {statistiques d'entraînement};
\end{tikzpicture}
\end{graphique}
\caption{Contenu d'un modèle enregistré, référence unique pour tout le document :
coefficients, six statistiques de préparation, ordre des matières et réponse
associée. Chaque élément est nécessaire à la prédiction.}
\label{fig:contenu-modele}
\end{figure}
\FloatBarrier

\section{Paramètres appris et paramètres fixés}

Les trois coefficients sortent du calcul ; tout le reste a été décidé avant lui,
et d'autres décisions auraient donné un autre modèle sur le même fichier.

\begin{table}[htbp]
\centering
\begin{tabular}{ll>{\raggedright\arraybackslash}p{54mm}}
\toprule
 & valeur ici & ce qui la fixe \\
\midrule
\multicolumn{3}{l}{\textbf{déterminé par la descente}}\\
\addlinespace[1pt]
coefficients $w$ & les trois valeurs ci-dessus
  & les notes et les maisons du fichier \\
\addlinespace[3pt]
\multicolumn{3}{l}{\textbf{posé avant la descente}}\\
\addlinespace[1pt]
matières retenues & \Herb, \Runes & le choix de deux matières qui se dessinent \\
remplissage des vides & la médiane & un arbitrage entre deux méthodes
  (ch.~\ref{ch:standardisation}) \\
échelle des notes & centrage-réduction & rendre les coefficients comparables \\
coefficients de départ & tous nuls & un point de départ quelconque convient \\
longueur du pas $\alpha$ & $1$ & un essai contrôlé sur la courbe de coût \\
critère d'arrêt & $\varepsilon$ et un plafond & le temps de calcul accordé \\
seuil de décision & $0{,}5$ & les deux erreurs comptées au même prix
  (ch.~\ref{ch:probabilite}) \\
\bottomrule
\end{tabular}
\caption{Le calcul produit trois nombres ; les sept décisions le précèdent.}
\label{tab:appris-pose}
\end{table}
\FloatBarrier

Ces décisions restent invisibles dans les trois coefficients enregistrés, et
doivent être consignées à côté d'eux pour qu'un tiers puisse refaire le
résultat.

\section{Classement d'une observation nouvelle}

Le calcul est celui de la figure~\ref{fig:carte}, la classe n'étant plus connue
mais déduite : standardiser les deux notes avec les statistiques d'entraînement,
calculer $z=w_0+w_1x_1+w_2x_2$, annoncer \Gryf si $z\geq0$. Cette comparaison à
zéro équivaut à comparer $p$ à $0{,}5$, la sigmoïde étant croissante ; le calcul
de $p$ est réservé à l'affichage.

Sur les $1\,600$ lignes du fichier, cette règle en classe correctement $1\,564$.
Les trente-six prédictions incorrectes relèvent de deux situations que la
distance à la frontière sépare.

La première est le recouvrement des deux groupes. Dans la bande où les deux
populations occupent la même région du plan, aucune droite ne peut les séparer
sans erreur, et déplacer la frontière ne fait que changer lesquelles.
Ces prédictions se reconnaissent à leur score, faible en valeur absolue, et à la
probabilité qui l'accompagne, voisine de $0{,}5$.

La seconde tient à des observations dont les notes contredisent l'étiquette.
Situées loin dans la région de l'autre groupe, elles reçoivent une prédiction
incorrecte assortie d'une probabilité élevée : ce sont les observations atypiques
dont la perte individuelle dépasse cent fois la
moyenne. La plus éloignée est incorrectement classée à une distance de
$2{,}27$ écarts-types de la frontière.

\begin{chiffres}[Erreurs par distance à la frontière]
\begin{center}
\begin{tabular}{lrl}
\toprule
distance & prédictions & \\
\midrule
moins de $0{,}5$ écart-type & $20$ & recouvrement des deux groupes \\
de $0{,}5$ à $1{,}0$        & $11$ & \\
plus de $1{,}0$             & $5$  & observations atypiques, jusqu'à $2{,}27$ \\
\midrule
\textbf{total}              & $\mathbf{36}$ & sur $1\,600$, soit $97{,}75\,\%$ de correctes \\
\bottomrule
\end{tabular}
\end{center}
La distance est comptée perpendiculairement à la frontière, en écarts-types, soit
$|z|/\|w\|$.
\end{chiffres}

Le taux de $97{,}75\,\%$ est l'\emph{exactitude sur l'échantillon
d'ajustement} : il sert de contrôle de cohérence, une valeur anormalement basse
signalant un défaut d'ajustement. L'exactitude hors échantillon, assortie de son
incertitude, est mesurée au chapitre~\ref{ch:predire}.

\section{Extension à dix notes et quatre maisons}

Passer de deux notes à dix ajoute huit coefficients et laisse les formules
identiques : la frontière devient un hyperplan de $\R^{10}$. Ce qui change ici
est la taille du fichier enregistré.

\begin{table}[htbp]
\centering
\begin{tabular}{lcccc}
\toprule
 & \multicolumn{2}{c}{\textbf{une maison}} & \multicolumn{2}{c}{\textbf{quatre maisons}} \\
\cmidrule(lr){2-3}\cmidrule(lr){4-5}
matières & coefficients & statistiques & coefficients & statistiques \\
\midrule
$2$  & $3$  & $6$  & $12$ & $6$  \\
$10$ & $11$ & $30$ & $44$ & $30$ \\
$n$  & $n+1$ & $3n$ & $4(n+1)$ & $3n$ \\
\bottomrule
\end{tabular}
\caption{Taille du modèle enregistré. Les coefficients se multiplient par le
nombre de maisons, les statistiques non : médiane, moyenne et écart-type sont
calculés une fois par matière et servent aux quatre modèles.}
\label{tab:taille}
\end{table}

Le passage à quatre maisons ajoute une décision : quatre modèles comme celui-ci
--- \Gryf contre le reste, puis \Huff contre le reste, et ainsi de suite ---
donnent quatre scores pour une même observation, et la classe retenue est celle
du plus grand (ch.~\ref{ch:classifieur}).

\begin{exercice}[Lecture d'un coefficient]
Un modèle ajusté sur les notes \emph{brutes} donnerait $w_2\approx+0{,}033$ pour
\Runes. Montrer que ce coefficient décrit la même frontière que $+3{,}469$, et
dire lequel des deux se compare à celui d'\Herb. Par quel facteur un écart-type
de plus en \Runes multiplie-t-il la cote, et pourquoi ce facteur ne dépend-il pas
de l'échelle retenue ?
\end{exercice}
